% ============================================================
% ECON 308 — International Trade
% Lecture 1 Quiz — INSTRUCTOR ANSWER KEY
% Fall 2026  |  Muhebullah Karimzada  |  UNBC
% ============================================================
% DO NOT DISTRIBUTE TO STUDENTS
% ============================================================
\documentclass[11pt]{article}
\usepackage[margin=1in,top=1.2in,bottom=1.2in]{geometry}
\usepackage[utf8]{inputenc}
\usepackage[T1]{fontenc}
\usepackage{lmodern}
\usepackage{amsmath,amssymb,bm}
\usepackage{booktabs,array,multirow,adjustbox}
\usepackage{enumitem}
\usepackage{xcolor}
\usepackage{fancyhdr}
\setlength{\headheight}{22pt}
\addtolength{\topmargin}{-10pt}
\usepackage{tcolorbox}
\usepackage{titlesec}
\usepackage{hyperref}
\tcbuselibrary{skins}
\hypersetup{colorlinks=false,hidelinks}

%% ── Colours ────────────────────────────────────────────────
\definecolor{UNBCDark}  {RGB}{0,74,37}
\definecolor{UNBCGold}  {RGB}{189,155,96}
\definecolor{SoftGray}  {RGB}{244,246,248}
\definecolor{CorrectGreen}{RGB}{0,120,60}
\definecolor{AccentRed}   {RGB}{192,57,43}
\definecolor{AccentBlue}  {RGB}{41,98,255}
\definecolor{AccentOrange}{RGB}{230,126,34}

%% ── Header / Footer ────────────────────────────────────────
\pagestyle{fancy}
\fancyhf{}
\renewcommand{\headrulewidth}{0.5pt}
\renewcommand{\headrule}{\hbox to\headwidth{\color{AccentRed}%
  \leaders\hrule height\headrulewidth\hfill}}
\lhead{\small\textbf{\color{UNBCDark}ECON 308 — Lecture 1 Quiz}%
       \;\textcolor{UNBCGold}{$|$}\;%
       \textcolor{AccentRed}{\bfseries INSTRUCTOR ANSWER KEY}}
\rhead{\small\color{UNBCDark}Fall 2026}
\cfoot{\small\thepage}

%% ── Section style ──────────────────────────────────────────
\titleformat{\section}[block]
  {\large\bfseries\color{UNBCDark}}{}
  {0em}{}
  [{\color{UNBCGold}\titlerule[0.8pt]}]
\titlespacing*{\section}{0pt}{1.2em}{0.6em}

\titleformat{\subsection}[runin]
  {\bfseries\color{UNBCDark}}{}
  {0em}{}[.\enspace]
\titlespacing*{\subsection}{0pt}{0.8em}{0.3em}

%% ── Custom boxes ───────────────────────────────────────────
\newtcolorbox{confidential}{
  colback=AccentRed!8, colframe=AccentRed,
  left=4mm, right=4mm, top=3mm, bottom=3mm,
  boxrule=0.8pt, arc=3mm}

\newtcolorbox{correctbox}{
  colback=CorrectGreen!10, colframe=CorrectGreen,
  left=3mm, right=3mm, top=1.5mm, bottom=1.5mm,
  boxrule=0.6pt, arc=2mm}

\newtcolorbox{solutionbox}{
  colback=AccentBlue!5, colframe=AccentBlue!60,
  title={\small\bfseries\color{UNBCDark}Solution},
  left=4mm, right=4mm, top=2mm, bottom=2mm,
  boxrule=0.6pt, arc=2mm}

\newtcolorbox{rubricbox}{
  colback=AccentOrange!8, colframe=AccentOrange,
  title={\small\bfseries Grading Rubric},
  left=3mm, right=3mm, top=2mm, bottom=2mm,
  boxrule=0.5pt, arc=2mm}

\newtcolorbox{databox}{
  colback=SoftGray, colframe=UNBCDark!40,
  left=4mm, right=4mm, top=3mm, bottom=3mm,
  boxrule=0.4pt, arc=2mm}

% ============================================================
\begin{document}
% ============================================================

%% ── Confidentiality banner ──────────────────────────────────
\begin{confidential}
\begin{center}
\large\bfseries\color{AccentRed}
INSTRUCTOR ANSWER KEY — DO NOT DISTRIBUTE TO STUDENTS
\end{center}
\end{confidential}

\vspace{0.5em}

%% ── Title block ─────────────────────────────────────────────
\begin{center}
{\LARGE\bfseries\color{UNBCDark}ECON 308 \;—\; International Trade}\\[0.5em]
{\large\bfseries Lecture 1 Quiz\quad\textcolor{UNBCGold}{$|$}\quad Fall 2026}\\[0.2em]
{\small University of Northern British Columbia\quad
 Instructor: M.\,Karimzada}
\end{center}

\vspace{0.3em}
{\color{UNBCGold}\hrule height 1.5pt}
\vspace{0.5em}

%% ── Quiz summary ─────────────────────────────────────────────
\begin{databox}
\small
\textbf{Quiz summary.}\quad
Source: Lec1.tex (Lec1.pdf), Slides/Lec1/, Fall 2026.\quad
Total: 20 marks.\quad
Time allowed: 20 minutes.\\[0.3em]
\textbf{Topics assessed:}
GDP expenditure identity and net exports (Q1, Part B(d));
gravity model equation and proportionality (Q2, Part B);
McCallum border effect (Q3);
Dutch Disease mechanism (Q5);
trade--macro transmission (Part B(d)).
\end{databox}

\vspace{0.6em}

%% ════════════════════════════════════════════════════════════
\section{Part A: Multiple Choice — Answer Key}
%% ════════════════════════════════════════════════════════════

\medskip
\begin{center}
\begin{tabular}{ccc}
\toprule
\textbf{Question} & \textbf{Correct Answer} & \textbf{Marks}\\
\midrule
1 & B & 2\\
2 & A & 2\\
3 & C & 2\\
4 & B & 2\\
5 & D & 2\\
\midrule
\textbf{Part A Total} & & \textbf{10}\\
\bottomrule
\end{tabular}
\end{center}

\bigskip

%% ── Q1 ───────────────────────────────────────────────────────
\noindent\textbf{Question 1.}\quad\textbf{Correct Answer: B}

\begin{correctbox}
\textbf{B) Trade \emph{deficit} of \$40 billion; this reduces GDP by \$40 billion.}
\end{correctbox}

\medskip
\noindent\textbf{Explanation.}\enspace
Net exports are $\text{NX} = X - M = 800 - 840 = -\$40$ billion. Since
$Y = C + I + G + \text{NX}$, a \$40 billion fall in NX directly reduces
$Y$ by \$40 billion (all else constant). This is a trade \emph{deficit}, not a
surplus (eliminating A and D). Option C is incorrect because exports and
imports are not subtracted separately; only the \emph{net} term $X - M$ enters
the identity.

\textit{Slide reference: ``The Trade--Macro Nexus'' (slide 2).}

\bigskip

%% ── Q2 ───────────────────────────────────────────────────────
\noindent\textbf{Question 2.}\quad\textbf{Correct Answer: A}

\begin{correctbox}
\textbf{A) Increase by exactly 20\%.}
\end{correctbox}

\medskip
\noindent\textbf{Explanation.}\enspace
With $\alpha = 1$, the gravity model is $T_{ij} = Y_i\cdot Y_j / D_{ij}$.
If $Y_i$ increases to $1.2\,Y_i$ (all else constant):
\[T_{ij}' = \frac{(1.2\,Y_i)\cdot Y_j}{D_{ij}} = 1.2 \cdot \frac{Y_i\cdot Y_j}{D_{ij}} = 1.2\,T_{ij}.\]
This is exactly a 20\% increase. Option B would require both $Y_i$ \emph{and}
$Y_j$ to grow, compounding the effect. Options C and D are economically
implausible and inconsistent with $\alpha = 1$.

\textit{Slide reference: ``Reading Gravity Coefficients'' (slide 10).}

\bigskip

%% ── Q3 ───────────────────────────────────────────────────────
\noindent\textbf{Question 3.}\quad\textbf{Correct Answer: C}

\begin{correctbox}
\textbf{C) National borders impose substantial non-tariff costs on trade, even under free trade agreements.}
\end{correctbox}

\medskip
\noindent\textbf{Explanation.}\enspace
McCallum's finding is known as the ``border effect'' or ``home bias puzzle.''
CUSFTA had already eliminated most tariffs, yet interprovincial trade was
approximately 20 times larger than Canada--U.S.\ trade at the same distance.
This demonstrates that tariffs are only part of the story: currency uncertainty,
different regulatory standards, separate legal systems, and customs formalities
all impose trade costs that persist even under FTAs. Anderson and van Wincoop
(2003) estimated the Canada--U.S.\ border acts as if it adds $\approx$75,000 km
of extra distance.

\textit{Slide reference: ``The Home Bias Puzzle'' (slide 16); ``Table 2-1'' (slide 18).}

\bigskip

%% ── Q4 ───────────────────────────────────────────────────────
\noindent\textbf{Question 4.}\quad\textbf{Correct Answer: B}

\begin{correctbox}
\textbf{B) World trade expanded from 1870 to 1914, then collapsed between 1914 and 1945 due to war and protectionism, before recovering and expanding again after World War II.}
\end{correctbox}

\medskip
\noindent\textbf{Explanation.}\enspace
Lecture 1 presents trade as unfolding in distinct waves. The \textbf{First Wave}
(1870--1914) was driven by falling transport costs (steamships, railways) and
the gold standard, pushing world trade to roughly 8\% of global GDP. The
\textbf{interwar collapse} (1914--1945) was caused by World War I, the Great
Depression, and protectionist policies such as the U.S.\ Smoot-Hawley tariff
(1930), which triggered widespread retaliation and drove world trade down by
about two-thirds. The \textbf{Second Wave} (post-1945) and \textbf{Third Wave}
(post-1990) saw trade recover and surpass pre-war levels, driven by GATT/WTO
liberalization, containerization, and the rise of global value chains.

\noindent\textbf{Why the distractors are wrong:}
\begin{itemize}[nosep,leftmargin=1.3em]
  \item \textbf{A)} Incorrect. Growth was \emph{not} continuous — the
        1914--1945 period saw a sharp collapse, not further expansion.
  \item \textbf{C)} Incorrect. Trade was already substantial before 1945;
        containerization (invented 1956) accelerated trade but was not its origin.
  \item \textbf{D)} Incorrect. Trade did \emph{not} recover immediately after
        World War I; the interwar period (1919--1945) saw sustained decline due
        to the Great Depression and protectionism.
\end{itemize}

\textit{Slide reference: ``Three Waves of Globalization'' (slides 19--22).}

\bigskip

%% ── Q5 ───────────────────────────────────────────────────────
\noindent\textbf{Question 5.}\quad\textbf{Correct Answer: D}

\begin{correctbox}
\textbf{D) A commodity boom raises foreign demand for the Canadian dollar, causing the CAD to appreciate, which makes Canadian manufactured exports more expensive for foreign buyers and reduces manufacturing competitiveness.}
\end{correctbox}

\medskip
\noindent\textbf{Explanation.}\enspace
Dutch Disease operates through the exchange rate channel, not directly through
factor or capital markets. The mechanism is: oil/commodity boom $\to$ increased
foreign demand for CAD (foreigners buy CAD to pay for Canadian oil) $\to$ CAD
appreciates $\to$ price of Canadian manufactured goods in foreign currency rises
$\to$ manufactured exports become less competitive $\to$ manufacturing output and
employment fall. During 2007--2012, CAD reached near-parity with USD; Ontario auto
and BC lumber exports fell sharply. Options A and B describe plausible economic
channels but are not the Dutch Disease mechanism as presented in Lecture 1.

\textit{Slide reference: ``Dutch Disease: When Resource Wealth Hurts Manufacturing'' (slide 31).}

\clearpage

%% ════════════════════════════════════════════════════════════
\section{Part B: Numerical Problem — Full Solution}
%% ════════════════════════════════════════════════════════════

\medskip
\noindent\textbf{Setup and Given Data}

\medskip
\begin{databox}
\small
Canada: $Y_C = \$2{,}400\text{B}$.
\quad Gravity model: $T_{ij} = A\cdot Y_i\cdot Y_j\,/\,D_{ij}$, $A=1$.

\medskip
\begin{tabular}{lccc}
\toprule
\textbf{Partner} & \textbf{GDP ($Y_j$, \$B)} & \textbf{Distance ($D$, km)}\\
\midrule
Australia   & 1{,}600  & 12{,}000 \\
Mexico      & 1{,}200  & 3{,}000  \\
United States & 25{,}000 & 1{,}000 \\
\bottomrule
\end{tabular}
\end{databox}

\bigskip

%% Part B(a) ───────────────────────────────────────────────────
\noindent\textbf{(a)}\quad\textbf{Compute $T_{\text{CA}}$, $T_{\text{CM}}$, $T_{\text{CU}}$.}
\quad(2 points)

\begin{solutionbox}
\textbf{Canada--Australia:}
\[T_{\text{CA}}
  = \frac{2{,}400 \times 1{,}600}{12{,}000}
  = \frac{3{,}840{,}000}{12{,}000}
  = \mathbf{320}\text{ (in \$billion-km units)}\]

\textbf{Canada--Mexico:}
\[T_{\text{CM}}
  = \frac{2{,}400 \times 1{,}200}{3{,}000}
  = \frac{2{,}880{,}000}{3{,}000}
  = \mathbf{960}\]

\textbf{Canada--United States:}
\[T_{\text{CU}}
  = \frac{2{,}400 \times 25{,}000}{1{,}000}
  = \frac{60{,}000{,}000}{1{,}000}
  = \mathbf{60{,}000}\]
\end{solutionbox}

\begin{rubricbox}
\small
\textbf{2 pts total.} Award \textbf{1 pt} for correct formulas applied to all three
pairs (partial credit for 2/3 correct). Award \textbf{1 pt} for correct arithmetic
throughout.
\end{rubricbox}

\bigskip

%% Part B(b) ───────────────────────────────────────────────────
\noindent\textbf{(b)}\quad\textbf{Ratio $T_{\text{CU}}/T_{\text{CA}}$ and decomposition.}
\quad(2 points)

\begin{solutionbox}
\textbf{Ratio:}
\[\frac{T_{\text{CU}}}{T_{\text{CA}}} = \frac{60{,}000}{320} = \mathbf{187.5}\]

\textbf{Decomposition into gravity factors:}
\[\frac{T_{\text{CU}}}{T_{\text{CA}}}
  = \underbrace{\frac{Y_{\text{US}}}{Y_{\text{Aus}}}}_{\text{Size factor}}
    \times
    \underbrace{\frac{D_{\text{CA}}}{D_{\text{CU}}}}_{\text{Proximity factor}}
  = \frac{25{,}000}{1{,}600}\times\frac{12{,}000}{1{,}000}
  = 15.625 \times 12
  = \mathbf{187.5\checkmark}\]

\textbf{Interpretation:}
\begin{itemize}[nosep,leftmargin=1.3em]
  \item \textbf{Size factor (15.6$\times$):} The U.S.\ economy is 15.6 times larger
        than Australia's. Both produce more goods to sell and can buy more imports,
        proportionally expanding bilateral trade.
  \item \textbf{Proximity factor (12$\times$):} Australia is 12 times farther from
        Canada than the U.S.\ Distance raises transport costs, information barriers,
        and time-in-transit costs, suppressing trade with Australia.
  \item \textbf{Combined:} The two factors multiply: Canada is predicted to trade
        \textbf{187.5 times more} with the U.S.\ than with Australia.
\end{itemize}
\end{solutionbox}

\begin{rubricbox}
\small
\textbf{2 pts total.} \textbf{0.5 pt} for correct ratio (187.5). \textbf{0.5 pt} for
correct algebraic decomposition. \textbf{1 pt} for clear economic interpretation of
both factors (size and distance). Accept equivalent verbal explanations that
correctly identify both factors.
\end{rubricbox}

\bigskip

%% Part B(c) ───────────────────────────────────────────────────
\noindent\textbf{(c)}\quad\textbf{CPTPP effect on Canada--Australia trade.}
\quad(2 points)

\begin{solutionbox}
\textbf{New effective distance:} $D_{\text{CA}}^{\text{CPTPP}} = 12{,}000 \times 0.75 = 9{,}000$ km.

\textbf{New predicted trade:}
\[T_{\text{CA}}^{\text{CPTPP}}
  = \frac{2{,}400 \times 1{,}600}{9{,}000}
  = \frac{3{,}840{,}000}{9{,}000}
  \approx \mathbf{426.7}\]

\textbf{Percentage increase:}
\[\%\Delta = \frac{T_{\text{CA}}^{\text{CPTPP}} - T_{\text{CA}}}{T_{\text{CA}}}
  \times 100
  = \frac{426.7 - 320}{320}\times 100
  = \frac{106.7}{320}\times 100
  \approx \mathbf{33.3\%}\]

\textit{Alternative derivation:} A 25\% reduction in distance multiplies $T$ by
$D_{\text{old}}/D_{\text{new}} = 12{,}000/9{,}000 = 4/3$,
so trade rises by $4/3 - 1 = 1/3 \approx 33.3\%$.

\medskip
\noindent\textbf{Economic interpretation:} CPTPP eliminates tariffs and reduces
regulatory barriers, cutting effective distance to Australia from 12,000 to 9,000 km.
This raises predicted Canada--Australia trade by one-third. Note, however, that even
after CPTPP, $T_{\text{CU}}/T_{\text{CA}}^{\text{CPTPP}} = 60{,}000/426.7 \approx 141$:
the US still dominates Canadian trade overwhelmingly, consistent with the gravity
model's prediction that proximity is difficult to overcome with policy alone.
\end{solutionbox}

\begin{rubricbox}
\small
\textbf{2 pts total.} \textbf{1 pt} for correct $T_{\text{CA}}^{\text{CPTPP}} \approx 426.7$.
\textbf{1 pt} for correct percentage increase (33.3\%; accept 33\% or $1/3$). Deduct
0.5 pt if student calculates a percentage decrease or uses 12,000 km incorrectly.
Bonus (no extra points): award full credit to students who note the policy
implication about US trade dominance persisting.
\end{rubricbox}

\bigskip

%% Part B(d) ───────────────────────────────────────────────────
\noindent\textbf{(d)}\quad\textbf{US recession shock to Canadian GDP.}
\quad(4 points)

\begin{solutionbox}
\textbf{Initial conditions:} $Y = \$2{,}400\text{B}$; $\text{NX} = +\$48\text{B}$;
$C + I + G = Y - \text{NX} = 2{,}400 - 48 = \$2{,}352\text{B}$.

\medskip
\textbf{Shock:} $\Delta X = -\$96\text{B}$; $\Delta M = 0$
$\;\Rightarrow\; \Delta\text{NX} = -\$96\text{B}$.

\medskip
\textbf{(i) New net exports:}
\[\text{NX}_{\text{new}} = +48 + (-96) = \mathbf{-\$48\text{ billion}}\]
\textit{(trade balance turns from surplus to deficit)}

\medskip
\textbf{(ii) New GDP:}
\[Y_{\text{new}} = (C + I + G) + \text{NX}_{\text{new}}
  = 2{,}352 + (-48) = \mathbf{\$2{,}304\text{ billion}}\]

\textit{Equivalently:} $\Delta Y = \Delta\text{NX} = -\$96\text{B}$,
so $Y_{\text{new}} = 2{,}400 - 96 = \$2{,}304\text{B}$. \checkmark

\medskip
\noindent\textbf{Summary:} The \$96B export fall causes a \textbf{\$96B direct
reduction in Canadian GDP} (a 4\% decline), shifting the trade balance from a
\$48B surplus to a \$48B deficit.
\end{solutionbox}

\begin{rubricbox}
\small
\textbf{4 pts total.} \textbf{2 pts} for (i): correct new NX = $-\$48$B (sign required).
\textbf{2 pts} for (ii): correct new GDP = \$2,304B. Accept either of the two approaches
shown. Common errors: (a) forgetting that $\Delta M=0$ so $\Delta\text{NX}=\Delta X$; (b) confusing the
sign convention for deficit vs. surplus.
\end{rubricbox}


\vspace{1em}
{\color{UNBCGold}\hrule height 0.6pt}

\section{Point Summary}

\begin{center}
\begin{tabular}{llcc}
\toprule
\textbf{Part} & \textbf{Question / Sub-part} & \textbf{Marks} & \textbf{Cumulative}\\
\midrule
A & Q1 — GDP identity and trade balance & 2 & 2\\
A & Q2 — Gravity model proportionality & 2 & 4\\
A & Q3 — McCallum border effect & 2 & 6\\
A & Q4 — Globalization waves (1870--present) & 2 & 8\\
A & Q5 — Dutch Disease mechanism & 2 & 10\\
\midrule
B & (a) Gravity calculations (3 pairs) & 2 & 12\\
B & (b) Ratio decomposition + interpretation & 2 & 14\\
B & (c) CPTPP effect + percentage change & 2 & 16\\
B & (d) Export shock to NX and GDP & 4 & 20\\
\midrule
  & \textbf{Total} & \textbf{20} & \\
\bottomrule
\end{tabular}
\end{center}

\vspace{0.8em}

\begin{databox}
\small\textbf{Suggested grade boundaries (20-mark quiz, 20 minutes):}\\[0.3em]
\begin{tabular}{@{}llll@{}}
A range: 18--20 & B range: 15--17 & C range: 12--14 & D: 10--11 \quad F: $<$10
\end{tabular}
\end{databox}

\vspace{1em}
{\color{UNBCGold}\hrule height 0.6pt}
\begin{center}
\small\itshape
--- End of Answer Key --- \quad ECON 308, Fall 2026, Lecture 1 Quiz
\end{center}

\end{document}
